\documentclass[11pt,a4paper]{article}
\usepackage[T1]{fontenc}
\usepackage{lmodern,amsmath,amssymb,microtype,booktabs}
\usepackage[margin=25mm]{geometry}
\usepackage[colorlinks=true,urlcolor=blue,linkcolor=blue]{hyperref}
\newcommand{\dd}{\mathop{}\!\mathrm d}
\newcommand{\Om}{\Omega}
\newcommand{\ii}{\mathrm i}
\title{Removable collisions in the dS$_3$ crossing kernel:\\merged series and a controlled heavy-mass factor}
\author{Alan Day \and Chengheng\thanks{AI research assistant (OpenAI).}}
\date{Working research note --- 31 October 2026}
\begin{document}\maketitle
\begin{abstract}
We reorganize the known LPRS scalar crossing formula to remove a
numerically unstable pair of Gamma poles. A termwise merged series has
a finite digamma-weighted limit on each real collision surface. In a
specified transfer window we prove an explicit absolute remainder bound
for one heavy-mass Omega factor, and a direct-sum tail bound that avoids
an uncertified asymptotic acceleration. A pointwise logarithmic
near-collision scaling follows for this factor. These results do not
establish the distributional limit of the full Racah transform or the
Bose free-observer limit. Numerical failures, including a prototype
precision bug, remain in the record.
\end{abstract}

\section{Physical scope and the removable surfaces}
The target remains the local $G_N=0$, radius-one model with proper
$SO^+(3,1)$, BD scalar observers and KG-normalized matter. Observer
labels exceed a fixed $s>0$; the two-creation sector uses the previously
defined half-period regulator and physical mass $\sqrt{1+\nu^2}$.
Here we study the scalar matrix element of the full crossing transform,
not a separately unitary scalar transform. No state, operator,
or physical spectral cutoff is changed.

In the LPRS formula \cite{lprs}, use
$h_i=(1+\ii\nu_i)/2$, $h=(1+\ii b)/2$,
$h'=(1+\ii d)/2$ and external labels $(a,q,c,p)$.
Both chiral weights are equal for the scalar principal series; they
are not complex conjugates. Each Omega factor in (3.38) is a sum
\[
 \Om(h,h';h_i)=\Gamma(z)A(z)+\Gamma(-z)B(z),\qquad
 z=h-h'+h_1+h_3-1.
\]
The full symbol contains $h$ and $1-h$, and $h'$ and $1-h'$.
For positive real labels the possible $z=0$ surfaces are
\[
 a+c=b+d,\qquad d=a+c+b,\qquad b=a+c+d .
\]
These are removable singularities of the two-term representation.
They are not physical poles at finite real mass.
The original Euler product integral (3.37) is locally integrable:
its absolute endpoint powers are $x^{-1/2}(1-x)^{-1/2}$,
multiplied at worst by logarithmic hypergeometric endpoint factors.
On compact parameter sets it gives continuity through these surfaces.

\section{Exact merger of the two series}
Fix $H$ and $h_i$, put $h'_0=H+h_1+h_3-1$ and $h'=h'_0-z$.
Define four upper parameters
\[
 (A_1,A_2,A_3,A_4)=
 (H+h_1+h_2-1,\ H+h_3+h_4-1,\ H+h_3-h_4,\ H+h_1-h_2).
\]
All have real part $1/2$ on our real-label slice.
Let
\[
 C(z)=\frac{\Gamma(1+z)\Gamma(1-z)\Gamma(2h'_0-2z)\Gamma(2H)}
 {\Gamma(A_1-z)\Gamma(A_2-z)\Gamma(A_3)\Gamma(A_4)} .
\]
Direct cancellation of Gamma factors in the two known $\!{}_4F_3$
series yields
\begin{align}
 \Om&=C(z)\sum_{n=0}^{\infty}\frac{G_n^A(z)-G_n^B(z)}{z},\label{merged}\\
 G_n^A(z)&=
 \frac{\prod_{j=1}^4\Gamma(n+A_j-z)}
 {\Gamma(n+2h'_0-2z)\Gamma(n+2H-z)
  \Gamma(n+1-z)\Gamma(n+1)},\\
 G_n^B(z)&=
 \frac{\prod_{j=1}^4\Gamma(n+A_j)}
 {\Gamma(n+2h'_0-z)\Gamma(n+2H)
  \Gamma(n+1+z)\Gamma(n+1)}.
\end{align}
The individual $G$'s have the same leading $n^{-2}$ behavior.
Their difference divided by $z$ is $2n^{-3}+O(n^{-4})$ for
fixed parameters, with the $z=0$ value interpreted continuously.
Thus the merged series converges absolutely and locally uniformly
near a real collision. This is an algebraic reorganization of an
existing formula, not a new definition of the physical kernel.

At $z=0$, the two regular prefactors $A(0)$ and $B(0)$ coincide.
Expanding $\Gamma(\pm z)$ gives
$A'(0)-B'(0)-2\gamma A(0)$.
Equivalently, termwise differentiation of \eqref{merged} gives
\begin{align}
 \Om_0&=\sum_{n=0}^{\infty}t_n\,q_n,\label{digamma}\\
 t_n&=\frac{\prod_{j=1}^4(A_j)_n}
 {(2h'_0)_n(2H)_n(n!)^2},\\
 q_n&=\psi(n+2h'_0)+\psi(n+2H)+2\psi(n+1)
       -\sum_{j=1}^4\psi(n+A_j).
\end{align}
The digamma bracket starts as $2/n+O(n^{-2})$.
The local uniform convergence justifies the parameter derivative.
No mass displacement is used to define the exact-surface value.

For computation away from zero, use
$G_n^B\operatorname{expm1}(\log G_n^A-\log G_n^B)/z$.
Logarithmic Gamma differences and successive ratios must retain their
branches continuously. Arbitrary precision still needs guard digits
when $|z|$ is tiny. Gamma-ratio/Bernoulli and Hurwitz-zeta acceleration
is available \cite{dlmf}, but its truncated large-parameter remainder
is not automatically controlled.

\section{A proved heavy-mass window}
Now select the colliding factor $H=(1-\ii b)/2$ with
\[
 c=M,\quad b=M+u,\quad a=M+u+v,\quad d=M+v,\qquad
 |u|+p\le1,\quad |v|+q\le1.
\]
Assume all physical observer labels remain above $s$.
Then $h'_0=(1+\ii d)/2$ and
\[
 (A_1,A_2,A_3,A_4)=
 \left(\frac{1+\ii(v+q)}2,\frac{1+\ii(p-u)}2,
       \frac{1-\ii(u+p)}2,\frac{1+\ii(v-q)}2\right).
\]
Write $A_j=1/2+\ii\alpha_j$; the window ensures $|\alpha_j|\le1/2$.
For $m=\min(b,d)>0$, the coefficient ratio obeys
\[
 \left|\frac{t_{n+1}}{t_n}\right|
 \le\left(\frac{n+1}{n+2}\right)^2
        \frac{(n+1)^2}{(n+1)^2+m^2}.
\]
To see this, bound each $|n+A_j|^2$ by $n^2+n+1/2$,
and use
$(x^2-x+1/2)(x+1)\le x^3$ for $x=n+1\ge1$.
Consequently
\begin{equation}
 |t_n|\le\frac1{(n+1)^2}
       \prod_{j=1}^n\frac{j^2}{j^2+m^2}
 \le\frac1{(n+1)^2(1+m^2)},\qquad n\ge1. \label{coeffbound}
\end{equation}

The digamma difference series gives the uniform bound
\[
 |q_n|\le B_M,\qquad
 B_M=2+\operatorname{arsinh}b+\operatorname{arsinh}d
             +4\log2\sum_j|A_j-1|.
\]
For the upper parameters, use
$\sum_{k\ge0}[(k+1/2)(k+1)]^{-1}=4\log2$.
For a denominator parameter $1+\ii y$, the difference from $\psi(n+1)$
is bounded by
$\sum_{j\ge1}|y|/[j\sqrt{j^2+y^2}]\le1+\operatorname{arsinh}|y|$.
Combining these estimates proves
\begin{align}
 \Om_0&=q_0+R_M,\qquad
 |R_M|\le\frac{(\zeta(2)-1)B_M}{1+m^2},\label{leading}\\
 q_0&=\psi(1-\ii b)+\psi(1+\ii d)+2\psi(1)-\sum_j\psi(A_j).
\end{align}
This is an explicit $O(\log M/M^2)$ absolute remainder for fixed
transfers in the stated window.
Standard digamma asymptotics then imply
$\Om_0=2\log M+2\psi(1)-\sum_j\psi(A_j)+O(M^{-1})$.

More usefully, summing \eqref{digamma} only through $n=N-1$ has
the analytic exact-arithmetic tail bound
\begin{equation}
 \left|\Om_0-\sum_{n=0}^{N-1}t_nq_n\right|
 \le \frac{B_M}{N}\prod_{j=1}^N\frac{j^2}{j^2+m^2}.
 \label{tail}
\end{equation}
Here $\sum_{n=N}^{\infty}(n+1)^{-2}\le1/N$.
Unlike a fitted zeta tail, \eqref{tail} is a proved bound on the
omitted series in this window. It does not bound floating-point error.

\section{Neighbouring masses and logarithmic scaling}
Let $d=M+v+r$ so that $z=-\ii r/2$.
Suppose both $|\operatorname{Im}A_j|\le1/2$ and
$|\operatorname{Im}(A_j-z)|\le1/2$.
The two unmerged hypergeometric series are $\!{}_4F_3$ with numerator
sets $A_j-z$ and $A_j$, and denominator sets
\[
 (2h',2H-z,1-z),\qquad (2h'_0-z,2H,1+z).
\]
Define
\begin{align*}
 a_0(z)&=\frac{\Gamma(2H)\Gamma(A_3-z)\Gamma(A_4-z)}
                {\Gamma(A_3)\Gamma(A_4)\Gamma(2H-z)},\\
 b_0(z)&=\frac{\Gamma(2h')\Gamma(A_1)\Gamma(A_2)}
                {\Gamma(A_1-z)\Gamma(A_2-z)\Gamma(2h'_0-z)},\\
 m_A&=\min(|M+v+r|,|M+u-r/2|),\\
 m_B&=\min(|M+v+r/2|,|M+u|).
\end{align*}
For $r\ne0$ and $m_A,m_B>0$, the same ratio proof gives
\begin{equation}
 \left|\Om-\Gamma(z)a_0-\Gamma(-z)b_0\right|
 \le(\zeta(2)-1)
 \left(\frac{|\Gamma(z)a_0|}{1+m_A^2}
      +\frac{|\Gamma(-z)b_0|}{1+m_B^2}\right). \label{off}
\end{equation}
Each hypergeometric tail has bound
$N^{-1}\prod_{j=1}^N j^2/(j^2+m_{A,B}^2)$,
multiplied by its corresponding Gamma prefactor.
These bounds grow as $1/|r|$ near zero;
use \eqref{tail} at exact collision.

For fixed $\lambda\ne0$, set $r=\lambda/\log M$.
Use the strict interior of the transfer window so the shifted
constraints eventually hold. Gamma-ratio asymptotics in the sectors
of $1\pm\ii M$ give
\[
 a_0(z)=(-\ii M)^z[1+O(z)+O(M^{-1})],\qquad
 b_0(z)=(\ii M)^{-z}[1+O(z)+O(M^{-1})].
\]
Since $\Gamma(\pm z)=\pm z^{-1}+O(1)$, the leading pair divided by
$2\log M$ tends to $\sin(\lambda/2)/(\lambda/2)$.
The error \eqref{off} is $O(\log M/M^2)$.
Together with \eqref{leading} at $\lambda=0$ this proves the
pointwise statement
\begin{equation}
 \frac{\Om\bigl(r=\lambda/\log M\bigr)}{2\log M}
 \longrightarrow \frac{\sin(\lambda/2)}{\lambda/2},
 \qquad \lambda\ \hbox{fixed}. \label{scaling}
\end{equation}
The right side is one at zero. We do not claim a uniform error in
$\lambda$, distributional convergence of the full $6j$, or a delta
replacement inside a physical mass integral.

\section{Numerical evidence and retained failures}
For $p=q=0.6$, $u=0.3$, $v=-0.2$, the complete physical kernel
on $a+c=b+d$ has the following pointwise reference values:
\begin{center}
\begin{tabular}{rrrrr}\toprule
$M$ & 2 & 4 & 8 & 16\\\midrule
$U_{00}$ & 0.96379723 & 1.15801422 & 1.36981862 & 1.58731944\\
\bottomrule\end{tabular}
\end{center}
A separate $M=32$ reference gives $1.80680785$.
These are point values, not a mass-integrated Bose observable.
The initial study also retains offsets $r=0.1,0.4$ at $M=2,4,8,16$.
No full-kernel limiting profile is inferred from these samples.

For the complete $M=32$ kernel, merely increasing arithmetic precision
does not repair an inadequate asymptotic tail. At 110 decimal digits,
the merged evaluator with $(N,K)=(512,24)$ has relative error about
$2.42\,10^{24}$ against the separate LPRS reference. The successive
settings $(1024,32)$, $(2048,40)$, $(4096,48)$ and $(8192,56)$ give
approximately $5.78\,10^6$, $1.74\,10^{-13}$, $5.39\,10^{-19}$ and
$8.60\,10^{-25}$. All attempts are retained. These are sample-specific
empirical errors, not a uniform accuracy claim for the algorithm.

For the single controlled Omega factor, the rigorous direct-sum
tail at $M=32$, $r=0$, $N=128$ is about $2.62\,10^{-39}$ in
exact arithmetic. At much larger $M$ this tail becomes far smaller
than the available rounding precision; it must not be reported as
the total numerical error. Computations at 70 and 100 decimal digits
are stored separately for the same series cutoff.

Convergence in \eqref{scaling} is slow: at $M=1024$ the complex scaled
value for $\lambda=0$ is approximately $1.31857+0.08411i$,
whose absolute difference from one is $0.32949$.
For $\lambda=1$ it is approximately $1.27753+0.09388i$,
with absolute difference $0.33222$ from $\sin(1/2)/(1/2)$.
The existence of the asymptotic limit does not imply these masses
are already close to it.


The first collision prototype accidentally used Python binary64
division in one Bernoulli coefficient before multiplication by an
arbitrary-precision value. At a low-mass sample it produced an
Omega bias about $2.82\,10^{-27}$ which was shared by precision
refinements. The old source and outputs remain unchanged.
The corrected merged module uses arbitrary-precision coefficients;
an independent parameter derivative reduces that sample discrepancy
to about $9.74\,10^{-58}$.

All three exact surfaces have been checked against that parameter
derivative, including matter-flip Hermiticity. Near-surface offsets
$10^{-8}$ and $10^{-25}$ and an independent Euler integral were also
checked. The nine diagnostics pass their stated tolerances; none is
an interval proof of the arbitrary-precision library.
An additional 100-digit comparison of the direct-sum analytic bounds
with complete hypergeometric or parameter-derivative evaluations is
still running at this checkpoint. It is not counted as a passed test.
No old physical wave-packet grid is overwritten.

\section{What remains before the Bose free limit}
The Bose target remains
$N_{\rm free}(t)=L_p(0)^2+|L_p(t)|^2$ and $D_{\rm free}(t)=0$
under an explicitly specified heavy-observer preparation protocol.
The present analytic bound covers only one of four Omega factors in
the complete LPRS expression. Other factors and shadow terms can
exponentially amplify individual absolute errors and cancel among
themselves. Their phases and physical normalization must be controlled
jointly. Pointwise kernel values cannot justify exchanging a mass
integral with a heavy-mass limit.

The next step is a uniform analysis of the full prefactored crossing
kernel against actual creation amplitudes, with mass tails and
cancellation-aware absolute bounds. A different integral
representation is preferable if the factorized representation cannot
supply such control. Complete Fock domains, finite gravity, and a
microscopic de Sitter holographic dual remain open.

\begin{thebibliography}{9}
\bibitem{lprs} J.~Liu, E.~Perlmutter, V.~Rosenhaus and D.~Simmons-Duffin,
\emph{$d$-dimensional SYK, AdS Loops, and $6j$ Symbols},
\href{https://arxiv.org/abs/1808.00612}{arXiv:1808.00612}, v3,
Eqs.~(3.35)--(3.38).
\bibitem{dlmf} NIST Digital Library of Mathematical Functions,
\href{https://dlmf.nist.gov/5.11}{Sec.~5.11}
(Gamma-ratio and digamma expansions), consulted 31 October 2026.
\end{thebibliography}
\end{document}
